\chapter*{现代化编写说明}
\addcontentsline{toc}{chapter}{现代化编写说明}
\markboth{现代化编写说明}{现代化编写说明}

本项目以庞特里亚金《常微分方程》全书为现代化重编目标，按原书章节逐步推进。
改写保留各章的正文、定理、命题、论证、例题、计算、讨论和图像；
现代改写主要处理三项内容：
把俄罗斯传统记号换成当前通用的向量与矩阵表达，把跨越多行的长句拆开，
并把原书留到后文的证明补成可以连续阅读的证明。

\noindent\textbf{一、章节编排与当前进度。}
当前版本已收录第一章、第二章及相关补证。
其中，第一章对应原书第 1--26 页，包含 \S1--\S6；
第二章对应原书第 27--99 页，包含 \S7--\S16。
第三章从原书 \S17 开始，将与后续章节按原书顺序继续编入。
下表记录当前已收录章节的对应关系，随编写进度扩充。

\begin{center}
\small
\setlength{\tabcolsep}{4pt}
\begin{tabularx}{\textwidth}{@{}p{2.5cm}p{3.1cm}X@{}}
\toprule
原书范围 & 本书位置 & 主要内容\\
\midrule
第 1--11 页 & 第一章 \S1--\S2 &
一阶方程、解与积分曲线、方向场、初等求积法\\
第 11--15 页 & 第一章 \S3 &
标准方程组、存在唯一性、延拓和线性整体存在\\
第 15--21 页 & 第一章 \S4 &
高阶方程与高阶系统化为一阶标准系统\\
第 21--25 页 & 第一章 \S5 & 复值微分方程\\
第 25--26 页 & 第一章 \S6 & 线性方程的基本性质\\
第 27--79 页 & 第二章 \S7--\S13 &
常系数方程、稳定性、非齐次方程、消去法、复振幅法和电路\\
第 71--99 页 & 第二章 \S14--\S16 &
标准线性系统、自治系统、相空间和相平面\\
\bottomrule
\end{tabularx}
\end{center}

\noindent\textbf{二、记号约定。}
向量与矩阵的符号统一使用粗正体，标量、向量分量和矩阵元素保持普通数学字体。
例如，$\vect{x}$ 是向量，$x_i$ 是其标量分量；$\mat{A}=(a_{ij})$ 是矩阵，
$a_{ij}$ 是其标量元素。标量函数通常写作 $y(t)$，向量解写作列向量
$\vect{x}(t)=(x_1(t),\ldots,x_n(t))^{\mathsf T}\in\mathbb F^n$。
带下标的整条向量解仍加粗，如 $\vect{x}_1(t)$，以区别于标量分量 $x_1(t)$。
矩阵从左侧作用于列向量。系统记为
\[
 \vect{x}'=\vect{f}(t,\vect{x}),\qquad \vect{f}:\Omega\to\mathbb F^n,
 \qquad (t,\vect{x})\in\Omega\subseteq\mathbb R\times\mathbb F^n.
\]
当需要强调对状态变量的微分时，记
\[
 D_{\vect{x}}\vect{f}(t,\vect{x})
 =\left(\frac{\partial f_i}{\partial x_j}(t,\vect{x})\right)_{i,j=1}^n.
\]
向量的时间导数一般写作 $\vect{x}'$、$\vect{x}''$ 或 $\vect{x}^{(k)}$。
在力学、电路、控制论和相空间等语境中，也保留原书风格的
$\dot{\vect{x}},\ddot{\vect{x}}$。两种写法含义相同，具体章节会优先采用与问题语义更贴近的一种。

常系数线性系统中的微分算子记为 $D=\dd/\dd t$。
为了区分算子与普通多项式变量，多项式写作 $L(s)$，
标量算子写作 $L(D)$，多项式矩阵及相应算子矩阵写作 $\mat{L}(s)$、$\mat{L}(D)$。
分量下标从 $1$ 开始；矩阵元素记作 $a_{ij}$。
复数的实部和虚部分别记作 $\operatorname{Re}$、$\operatorname{Im}$，
共轭记作 $\overline z$。

\noindent\textbf{三、改写原则。}
\begin{enumerate}[label=(\arabic*)]
  \item 保留原书的节号、定理号、命题次序和例题顺序，便于逐页对读。
  \item 只有数学对象完全相同时才更换记号；不把原书的论证替换成更短但不同的证明。
  \item 高阶标量方程写成状态向量系统，线性系统写成矩阵方程，
        多个未知函数的分量统一使用列向量。
  \item 原书的几何解释和相图保留。图题写清坐标轴、特征方向、
        轨线方向、平衡位置和参数条件。
  \item 原书本身省略的证明，在章末以“补证”给出并明确标注；
        不把补证冒充成原书正文。明显笔误另以“编者注”说明。
  \item 长句可以拆分，省略语可以补足，但不删去原书的条件、
        反例、退化情形和适用边界。
\end{enumerate}

\noindent\textbf{四、证明的编排。}
原书第一章把定理 1--3 的证明安排在第四章。
本书把所需的 Picard 逐次逼近、唯一性、最大延拓和线性系统整体存在
集中放在第一章正文之后，使第一章可以独立阅读。
原书第二章的 Hurwitz 判据只陈述而不证明；本书在第二章正文之后
补上基于 Routh 降阶的证明。补证采用与原书相同的定理条件，
并说明它与原书论证的关系。
